\subsection{同构的扩张}

命题8. --- 设 \(\Omega\) 是域 K 的一个代数闭扩张，E 和 F 是 \(\Omega\) 的两个子扩张， \(u\) 是 E 到 F 上的一个 K-同构。为使一个 K-自同构 \(v\) （ \(\Omega\) 的、且延拓 \(u\) ）存在，充要条件是 \(\Omega\) 在 E 与 F 上具有相同的超越次数。

该条件显然是必要的。

现设 \(\Omega\) 在 E 和 F 上具有相同的超越次数，并选取 \(\Omega\) 在 E 上的一个超越基 B 以及 \(\Omega\) 在 F 上的一个超越基 C。由于 B 与 C 等势，命题2（V，第106页）表明 \(u\) 可扩张为 E(B) 到 F(C) 上的一个 K-同构 \(u'\) 。由于 \(\Omega\) 是 E(B) 与 F(C) 的代数闭包，V 第23页的推论表明 \(u'\) 可扩张为一个自同构 \(v\) （ \(\Omega\) 的）。

推论1. --- 设 \(\Omega\) 是域 K 的一个代数闭扩张，E 是 \(\Omega\) 的一个子扩张。 E 的每个 K-自同构都可扩张为 \(\Omega\) 的一个 K-自同构。

推论2. --- 设 \(\Omega\) 是域 \(K\) 的一个代数闭扩张， \(E\) 与 \(F\) 是 \(\Omega\) 的两个子扩张， \(u\) 是 \(E\) 到 \(F\) 上的一个 K-同构。若 \(E\) 在 \(K\) 上的超越次数有限（特别地，若 \(E\) 在 \(K\) 上是代数的），则存在 \(\Omega\) 的一个延拓 \(u\) 的 K-自同构。

我们用 n、 \(d(\mathrm{E})\) 与 \(d(\mathrm{F})\) 分别表示 E 在 K 上的、 \(\Omega\) 在 E 上的以及 \(\Omega\) 在 F 上的超越次数。 K-同构 u 的存在性表明 F 在 K 上的超越次数等于 n。由定理4的推论， \(\Omega\) 在 K 上的超越次数等于 \(d(\mathrm{E}) + n\) ，也等于 \(d(\mathrm{F}) + n\) 。因此（E，III，第28页，命题8）我们有 \(d(\mathrm{E}) = d(\mathrm{F})\) ，从而可以应用命题8。

命题9. --- 设 K 是一个域， \(\Omega\) 是 K 的一个代数闭扩张。假设 \(\Omega\) 在 K 上不是代数的；则 \(\Omega\) 中在 K 上超越的元素所成的集合是无限的。此外，若 \(x\) 与 \(y\) 是 \(\Omega\) 的两个在 K 上超越的元素，则存在一个自同构 \(u\) （ \(\Omega\) 的、在 K 上的），使得 \(u(x) = y\) 。

由于 \(\Omega\) 在 \(\mathbf{K}\) 上不是代数的，存在一个元素 \(x\) （ \(\Omega\) 的、在 \(\mathbf{K}\) 上超越的）；于是元素 \(x^n\) （ \(n \in \mathbf{N}\) ）两两不同且在 \(\mathbf{K}\) 上超越。设 \(x\) 与 \(y\) 不同且在 \(\mathbf{K}\) 上超越；由命题2（V，第106页），存在一个 K-同构 \(\tilde{u}\) ，从 \(\mathbf{K}(x)\) 到 \(\mathbf{K}(y)\) 上，使得 \(\tilde{u}(x) = y\) ；由于 \(\mathbf{K}(x)\) 在 \(\mathbf{K}\) 上的超越次数为 1，命题8的推论2表明 \(\tilde{u}\) 可扩张为一个 K-自同构 \(u\) （ \(\Omega\) 的）。

命题10. --- 设 K 是一个域，Ω 是 K 的一个代数闭扩张，G 是 Ω 的 K-自同构群。设 \(x \in \Omega\) 。

\begin{enumerate}
\def\labelenumi{\alph{enumi})}
\item
  为使 \(x\) 在 \(K\) 上是代数的，充要条件是：元素 \(u(x)\) （当 \(u\) 取遍 \(G\) ）所成的集合是有限的，
\item
  为使 \(x\) 是 \(p\) -根式的（相对于 \(\mathbf{K}\) ），充要条件是 \(u(x) = x\) （对所有 \(u \in \mathbf{G}\) ）成立。
\end{enumerate}

特别地，若 K 是完全域，则群 G 的不变量集合等于 K。首先设 x 是超越的。由命题9， \(\Omega\) 中在 K 上超越的元素所成的集合 T 是无限的，且对每个 \(y \in T\) ，存在 \(u \in G\) 使得 \(u(x) = y\) 。因此，当 u 取遍 G 时元素 \(u(x)\) 所成的集合是无限的。

现设 \(x\) 在 \(\mathbf{K}\) 上是代数的，并用 \(f\) 表示它在 \(\mathbf{K}\) 上的极小多项式； \(f\) 在 \(\Omega\) 中的根所成的集合是有限的，且对每个 \(u \in \mathbf{G}\) 我们有 \(f(u(x)) = u(f(x)) = 0\) 。因此，元素 \(u(x)\) （当 \(u\) 取遍 \(\mathbf{G}\) ）所成的集合是有限的。这就证明了 \(a\) 。

设 L 为 \(\Omega\) 中满足 \(u(y)=y\) （对所有 \(u\in G\) ）的元素 y 组成的集合，并设 \(\bar{K}\) 为 K 在 \(\Omega\) 中的相对代数闭包。由上所述，L 是 \(\bar{K}\) 在 K 上的一个子扩张。此外（命题8的推论1）， \(\bar{K}\) 的每个 K-自同构都是 G 的某个元素在 \(\bar{K}\) 上的限制。命题10的断言 b) 现由 V 第53页的推论3得出。

